\documentclass[8pt]{extarticle}
\usepackage[utf8]{inputenc}
\usepackage[spanish]{babel}
\usepackage{amsmath, amssymb}
\usepackage{graphicx}
\usepackage{multicol}
\usepackage{tikz}
\usetikzlibrary{babel}
\usepackage{xskak}
\usepackage{chessboard}
\usepackage{tcolorbox} 
\usepackage{fontawesome5} 
\usepackage{bbding}
\usepackage{fourier-orns}
\usepackage{mathrsfs}
\usepackage[shortlabels]{enumitem}
\usepackage{enumerate} 
\usepackage{booktabs}
\usepackage{array}

% Tamaño de cada mini-página (1/8 de hoja oficio de 21.6 x 34 cm)
\usepackage[paperwidth=8.5cm, paperheight=10.8cm, margin=0.6cm]{geometry}
\pagestyle{empty} % Quita los números de página grandes para ahorrar espacio

\setlength{\columnsep}{1pt}

\begin{document}

% --- PÁGINA 1: PORTADA ---
\begin{center}
    \vspace*{-0.3cm}
   {\fontsize{26}{28}\selectfont 
    \textbf{El Semillero}\\[-4pt]
    {\fontsize{20}{22}\selectfont \textbf{de}}\\[-3pt]
    \textbf{Grothendieck}\par}

\vspace{0.4cm}

{\large\textit{``Récoltes et semailles''}}

\vspace{0.5cm}

\begin{center}
\begin{tikzpicture}[scale=0.8]
    % Marco invisible simétrico alrededor del centro del polígono (x = 0.05)
    \useasboundingbox (-2.45,-2.1) rectangle (2.55,2.3);

    % Definición de las raíces de P(z)
    \coordinate (z1) at (-2.3, 0.8);
    \coordinate (z2) at (-0.8, 1.8);
    \coordinate (z3) at (2.1, 1.3);
    \coordinate (z4) at (1.4, -0.9); % <-- MOVIDO A LA IZQUIERDA (De 2.4 a 1.4)
    \coordinate (z5) at (-1.5, -1.6);

    % Envolvente convexa en gris claro
    \filldraw[fill=gray!15, draw=black, thick] 
        (z1) -- (z2) -- (z3) -- (z4) -- (z5) -- cycle;

    % Ejes del plano complejo (simétricos respecto al polígono)
    \draw[->, thin, gray!70] (-2.8,0) -- (2.9,0) node[right, black] {\tiny $\mathrm{Re}$};
    \draw[->, thin, gray!70] (0,-2.0) -- (0,2.2) node[above, black] {\tiny $\mathrm{Im}$};

    % Raíces de P(z) (círculos negros sólidos)
    \foreach \p/\pos/\nombre in {
        z1/left/$z_1$, 
        z2/above/$z_2$, 
        z3/right/$z_3$, 
        z4/below/$z_4$, % <-- CAMBIADO A BELOW para que no se encime con la línea
        z5/below/$z_5$
    }{
        \fill[black] (\p) circle (2.2pt) node[\pos, black] {\footnotesize \nombre};
    }

    % Raíz interior extra de P(z)
    \coordinate (z6) at (-0.6, 0.2);
    \fill[black] (z6) circle (2.2pt) node[left, black] {\footnotesize $z_6$};

    % Raíces de la derivada P'(z) (cuadrados blancos)
    \coordinate (w1) at (-1.1, 0.6);
    \coordinate (w2) at (0.5, 1.0);
    \coordinate (w3) at (1.1, -0.1);
    \coordinate (w4) at (-0.3, -0.7);
    \coordinate (w5) at (0.1, 0.2);

    \foreach \w in {w1, w2, w3, w4, w5}{
        \filldraw[fill=white, draw=black, thick] (\w) ++(-1.8pt,-1.8pt) rectangle ++(3.6pt,3.6pt);
    }

    % Etiquetas
    \node[anchor=west] at (1, 2.0) {\footnotesize $\mathrm{Conv}(\{z_{1},z_{2},\dots ,z_{6}\})$};
    \node[black, above right] at (w2) {\small $w$};
\end{tikzpicture}

\end{center}


\vspace{0.4cm}

\noindent
    {\fontfamily{ppl}\selectfont\footnotesize\bfseries
\parbox[c]{1.7cm}{\centering\textbf{Esenciales del \\ Estudiante de Licenciatura}}%
\hfill
\parbox[c]{1.7cm}{\centering\textbf{Jardín de \\Espacios Topológicos}}%
\hfill
\parbox[c]{1.7cm}{\centering\textbf{El Teorema de \\Gauss-Lucas}}%
\hfill
\parbox[c]{1.7cm}{\centering\textbf{Matrices Simétricas y\\  Antisimétricas}}%
}
    \vspace{0.2cm}
    \hrule
    \vspace{0.2cm}
    {\footnotesize Número 1 -- Noviembre 2026}
\end{center}
\vspace{0.5cm}


\newpage

% --- PÁGINA 2 ---
\begin{center}
    {\large \textbf{ $\S$ Recomendación}} {$~~$\small\faCoffee}
\end{center}

\vspace{0.14cm}

\begin{center}
    {\large\textit{\textbf{Essential Mathematics for Undergraduates:}}\\ \vspace{0.3cm}
    {\small A Guided Approach to Algebra, Geometry,\\ Topology and Analysis.}}\\
    \vspace{0.3cm}
    {\small \textbf{Simon G. Chiossi}}
    
\end{center}

\vspace{0.14cm}

{\small \fontfamily{qpl}\selectfont Nuestra primera recomendación...}
\vspace{0.2cm}

\noindent
{\small \fontfamily{jkp}\selectfont \textbf{Theorem 3.10$~~$(Gödel)} $~~$\textit{The negation of the axiom of choice $\neg \mathrm{AC}$ is not deducible in $\mathrm{ZF}$.}} 
\vspace{0.2cm}

{\small \fontfamily{qpl}\selectfont Por supuesto... $\dagger$}



\newpage

% --- PÁGINA 3 ---

\begin{center}
    {\large \textbf{ $\S$ Jardín de Estructuras}} {$~$\small\FourClowerSolid }
\end{center}
\vspace{0.14cm}

\begin{center}
    {\decothreeleft\large\textit{ Espacios Topológicos }\decothreeleft}
\end{center}

\vspace{0.15cm}

{\fontsize{7}{5} \fontfamily{jkp}\selectfont  Recordemos que si $X$ es un conjunto... 
Entre algunos ejemplos tenemos:}


\begin{center}
\footnotesize \fontfamily{jkp}\selectfont 
\setlength{\tabcolsep}{4pt}
\begin{tabular}{@{} c c m{2.55cm} @{}}
\toprule
\textbf{Espacio} & \textbf{Notación} & \multicolumn{1}{c}{\textbf{Descripción}} \\
\midrule
Discreto & $(X, \mathscr{P}(X))$ & \raggedright\arraybackslash\scriptsize Metrizable y $0$-dimensional. \\[4pt]
Indiscreto & $(X, \{\varnothing, X\})$ & \raggedright\arraybackslash\scriptsize Pseudometrizable y compacto. \\
\bottomrule
\end{tabular}
\end{center}


\newpage

% --- PÁGINA 4 ---

\begin{center}
    {\large \textbf{ $\S$ El Teorema de Gauss-Lucas}} {$~$\small\faPenNib}
\end{center}

{\fontsize{7}{5} \fontfamily{jkp}\selectfont
Bajo la influencia... 
}
\newpage

% --- PÁGINA 5 ---
{\fontsize{7}{5} \fontfamily{jkp}\selectfont
Por otro lado... 
}
\vspace{0.2cm}

\noindent
{\small \fontfamily{jkp}\selectfont \textbf{Teorema$~$(Gauss-Lucas)} $~~$\textit{Sea $p\in \mathbb{C}[x]$ tal que $p\notin \mathbb{C}$. Entonces, para cada $w\in \mathbb{C}$ tal que $p'(w)=0$,
$$w\in \mathrm{Conv}(\{z\in \mathbb{C}\mid p(z)=0\}).$$
}} 
\vspace{0.2cm}
{\fontsize{7}{5} \fontfamily{jkp}\selectfont Véase... 

{\footnotesize \fontfamily{jkp}\selectfont \textbf{Referencia.}
Bak, J., y Newman, D. J. (2010). \textit{Complex analysis} (3.ª ed.). Springer.} 


}
\newpage

% --- PÁGINA 6 ---

\begin{center}
    {\large \textbf{ $\S$ Ejercicio para el Lector}} {$~$\small\faGlasses}
\end{center}

\begin{center}
    Álgebra Lineal $~~~~~~~~$ Dificultad: $\bullet \bullet \circ \circ\circ $
\end{center}

{\fontsize{7}{5} \fontfamily{jkp}\selectfont Dado $F$ un campo...
}
\newpage

% --- PÁGINA 7 ---
{\fontsize{7}{5} \fontfamily{jkp}\selectfont Observa además...
}
\vspace{0.2cm}

\noindent
{\small \fontfamily{jkp}\selectfont \textbf{Ejercicio.} $~$\textit{Sea $F$ un campo tal que $\mathrm{char}(F)\neq 2$. Si denotamos
$$\mathcal{A}_{n}(F):=\{A\in \mathcal{M}_{n}(F)\mid A^{t}=-A\}$$
entonces $\mathcal{M}_{n}(F)=\mathcal{S}_{n}(F)\oplus \mathcal{A}_{n}(F)$.}
\vspace{0.1cm}

\noindent
\textbf{Hint.} \textit{Toma $A\in \mathcal{M}_{n}(F)$. ¿De que manera puedes construir una matriz simétrica unicamente en términos de $A$?} 
}



\newpage

% --- PÁGINA 8: CONTRAPORTADA ---


\begin{figure}[h]
    \raggedleft
    \vspace{1.3cm} % <-- Baja la imagen la distancia que quieras
    % \includegraphics[width=0.5\linewidth]{Rama de cerezo en tinta gris minimalista.png}\hspace*{-1.8cm}
    \label{fig:placeholder}
    \vspace{-1.3cm} % <-- ¡TRUCO! Cancela el espacio para que lo de abajo NO se mueva
\end{figure}


\noindent
\raisebox{-4.7cm}[0pt][0pt]{%
    % \hspace*{-2cm}\reflectbox{\includegraphics[width=0.5\linewidth]{Rama de cerezo en tinta gris minimalista.png}}%
}

\vspace{-3.3cm}


\begin{center}

\newchessgame
\hidemoves{1.e4 e5}\par
\storegame{game1}
\chessboard[
label=false,         % <-- OCULTA LOS NÚMEROS (1-8)
  boardfontsize=7pt,
  setfen=k1b4r/pp6/3B4/P1n2p1q/2N5/6P1/5P2/R1R2BK1 w - - 0 20
]
\end{center}
\vspace{-0.4cm}
\begin{center}
    {\footnotesize\textbf{Mate en 7}}
\end{center}


\vspace{0.1cm}

\begin{tcolorbox}[
    colback=white,          % Fondo blanco
    colframe=black,   
     boxrule=0.2mm,    % Borde negro
    arc=0mm,                % Esquinas rectas (sin redondeado)
    width=0.7\textwidth,    % Ajusta el ancho del rectángulo según lo necesites
    halign=center,          % Centra el contenido dentro de la caja
    valign=center,
    center                  % Centra la caja misma en la página
]
    {\large\textbf{¿Te interesa colaborar?}}
    
    \vspace{0.2cm}
    \footnotesize
    
    {\fontfamily{ppl}\selectfont Editor}
    
    \vspace{0.1cm}
    {\textit{Yves Stanislas Ramírez Jiménez}}
    
    \vspace{0.1cm}
    {\fontfamily{ppl}\selectfont Instagram: math\_digital.garden}
    
    \vspace{0.1cm}
    {\fontfamily{ppl}\selectfont Web: Math Digital Garden}
\end{tcolorbox}

\vspace{0.1cm}

\begin{center}
    {\small Más matemáticas...}
\end{center}

\vspace{-0.1cm}

\begin{center}
       {\small \selectfont 
    \textbf{El Semillero}\\
    \textbf{de}\\
    \textbf{Grothendieck}\par}
\end{center}



\end{document}
